% Frequency converter: conversion transconductance (heptode-conversion): the set-up a person would wire, semantically the deck's.
%   VOSC osc 0 SIN(0 <Vosc> <f_osc>), CO osc g1 …, RG1 g1 0 20k   the oscillator grid, grid-leak biased   -> source (set), ammeter I_g1 (measured)
%   VSIG g3 0 DC <vg3> SIN(…)                                      the signal grid                         -> source (set)
%   VS g2 0 <vg2>, VA a 0 <va>, XT a g3 g2 g1 0 DUT                the screen and the anode                -> ammeter I_a (measured)
%   .tran; the anode current's component at |f_osc - f_sig| over the signal amplitude
\documentclass[tikz, border=10pt]{standalone}
\usepackage{bench}
\begin{document}
\begin{tikzpicture}
\benchtitle{0,7.0}{Frequency converter: conversion transconductance}
\def\gnd{1.0}
% the pentode
\node[pentode, scale=1.6] (T) at (6.4,3.0) {};
\node[anchor=west, font=\small, align=left] at ($(T.anode)!0.62!(T.cathode)+(0.95,0)$) {DUT\\(heptode)};
% the cathode to ground
\draw (T.cathode) -- (T.cathode |- 0,\gnd) node[ground]{};
\nodename{anchor=west}{$(T.cathode)+(0.1,-0.35)$}{0}
% the ammeter in the anode, up to the rail
\coordinate (top) at ($(T.anode)+(0,1.6)$);
\ammeter{T.anode}{top}{$I_a$}
\nodename{anchor=west}{$(T.anode)+(0.1,0.25)$}{a}
% the anode supply at the right, + up, - to ground
\coordinate (vtop) at ($(top)+(4.2,0)$);
\draw (top) -- (vtop);
\draw[set] (vtop) to[V, l={$V_a$}] (vtop |- 0,\gnd) node[ground, black]{};
% the signal grid and the screen as labelled leads (their sources in the definitions)
\draw (T.suppressor) -- ++(-0.3,0) -- ++(0,1.3) node[above left, set, font=\small, inner sep=1pt] {$g_3$: $V_{g3} + v_{sig}$};
\draw (T.screen) -- ++(-1.6,0) -- ++(0,1.0) node[above, set, font=\small, inner sep=1pt] {$g_2$: $V_{g2}$};
% the oscillator grid: its 20 k leak to ground, the ammeter, the coupling capacitor and the oscillator
\coordinate (g) at (T.control);
\draw (g) -- ++(-0.6,0) coordinate (gj);
\draw (gj) to[R, l={20 k\ohm}] (gj |- 0,\gnd) node[ground]{};
\coordinate (gam) at ($(gj)+(-1.7,0)$);
\ammeterh{gam}{gj}{$I_{g1}$}
\draw (gam) to[C, l_={$C$}] ++(-1.4,0) coordinate (osc);
\draw[set] (osc) to[sV, l_={$v_{osc}$}] (osc |- 0,\gnd) node[ground, black]{};
\nodename{anchor=south}{$(g)+(-0.3,0.02)$}{g1}
% the definitions
\definitions{0,-0.1}{%
  $V_a = V_{g2} = 250$ V (the 6BE6 sheet's columns);\quad the oscillator set for $I_{g1} = 0.5$ mA (10 V rms into 20 k\ohm);\quad $V_{g3} = -1.5$ to $-30$ V\\[2pt]
  read the conversion transconductance $g_c$ against $V_{g3}$ and against the oscillator grid current;\quad only for a model with a separate signal grid (Reefman's heptode)%ENDDEFS
}
\end{tikzpicture}
\end{document}
